\documentclass[../hippo.tex]{subfiles}
\begin{document}

\section{Experiment Details and Additional Results}
\label{sec:experiment-details}


\subsection{Model Architecture Details}
\label{subsec:model_architectures}

Given inputs $x_t$ or features thereof $f(x_t)$ in any model, the HiPPO framework can be used to memorize the history of features $f_t$ through time.
As the discretized HiPPO dynamics form a linear recurrent update similar in style to RNNs (e.g., \cref{thm:legs}),
we focus on these models in our experiments.

Thus, given any RNN update function $h_{t} = \tau(h_{t-1}, x_t)$, we simply replace the previous hidden state with a projected version of its entire history.
\begin{figure}[!ht]
  \begin{minipage}{.5\linewidth}%
    \centering
    \includegraphics[width=\linewidth]{figures/rnncell2.pdf}
  \end{minipage}
  \begin{minipage}{.5\linewidth}
    \begin{equation}%
      \label{eq:architecture}
      \begin{aligned}
        h_t \in \R^d &= \tau(h_{t-1}, [c_{t-1}, x_t]) \\
        f_t \in \R^1 &= \mathcal{L}_f(h_{t}) \\
        c_t \in \R^N &= \hippo_t(f) \\
        &= A_t c_{t-1} + B_t f_t
      \end{aligned}
    \end{equation}
  \end{minipage}
  \caption{The simple RNN model we use HiPPO with, and associated update equations. $\mathcal{L}_\square$ is a parametrized linear function, $\tau$ is any RNN update function, and $[\cdot]$ denotes concatenation.
    $\textcolor{magenta}{\hippo}$ is the HiPPO memory operator which orthogonalizes the history of the $f_t$ features up to time $t$. $A_t, B_t$ are fixed matrices depending on the chosen measure .
    $N$ and $d$ represent the approximation order and hidden state size, respectively.
  }
  \label{fig:cell-equation}
\end{figure}
Equations~\eqref{eq:architecture} lists the explicit update equations and Figure~\ref{fig:cell-equation} illustrates the model.
In our experiments, we choose a basic gated RNN update
\begin{align*}
  \tau(h, x) = (1-g(h, x)) \circ h + g(h, x) \circ \tanh(\mathcal{L}_\tau(h, x)), \qquad g(h, x) = \sigma(\mathcal{L}_g(h, x)).
\end{align*}

\paragraph{Methods and Baselines}

We consider the following instantiations of our framework HiPPO.

\textbf{HiPPO-}\textbf{LegT}, \textbf{LagT}, and \textbf{LegS}, use the translated Legendre, and tilted Laguerre, and scaled Legendre measure families with update dynamics \eqref{eq:translated-legendre-dynamics}, \eqref{eq:laguerre-dynamics}, and \eqref{eq:scaled-legendre-dynamics}.
As mentioned, LegT has an additional hyperparameter $\theta$, which should be set to the timescale of the data if known a priori.
We attempt to set it equal to its ideal value (the length of the sequences) in every task,
and also consider $\theta$ values that are too large and small
to illustrate the effect of this hyperparameter.

Our derivations in \cref{sec:derivation-legt,sec:derivation-lagt,sec:derivation-legs,sec:derivation-fourier,sec:derivation-chebyshev} show that there is a large variety of update equations that can arise from the HiPPO framework---for example, the tilted generalized Laguerre polynomials lead to an entire family governed by two free parameters (\cref{sec:derivation-lagt})---many of which lead to linear dynamics of the form $\ddt c(t) = -A c(t) + B f(t)$ for various $A, B$.
Given that many different update dynamics lead to such dynamical systems that give sensible results,
we additionally consider the \textbf{HiPPO-Rand} baseline that uses random $A$ and $B$ matrices (normalized appropriately) in its dynamics.

We additionally compare against the following standard RNN baselines.
The \textbf{RNN} is a vanilla RNN.
The \textbf{MGU} is a minimal gated architecture, equivalent to a \textbf{GRU} without the reset gate. The HiPPO architecture we use is simply the MGU with an additional $\hippo$ intermediate layer.
The \textbf{LSTM} is the most well-known and popular RNN architecture, which is a more sophisticated gated RNN.
The \textbf{expRNN}~\citep{lezcano2019cheap} is the state-of-the-art representative of the \emph{orthogonal RNN} family of models designed for long-term dependencies~\citep{arjovsky2016unitary}.
The \textbf{LMU} is the exact same model as in \citet{voelker2019legendre}; it is equivalent to HiPPO-LegT with a different RNN architecture.

All methods have the same hidden size in our experiments.
In particular, for simplicity and to reduce hyperparameters, HiPPO variants tie the memory size $N$ to the hidden state dimension $d$.
The hyperparameter $N$ and $d$ is also referred to as the number of hidden units.

\paragraph{Model and Architecture Comparisons}
The model~\eqref{eq:architecture} we use is a simple RNN that bears similarity to the classical LSTM and the original LMU cell.
In comparison to the LSTM, HiPPO can be seen as a variant where the memory $m_t$ plays the role of the LSTM's hidden state and $h_t$ plays the role of the LSTM's gated cell state, with equal dimensionalities.
HiPPO updates $m_t$ using the fixed $A$ transition matrix instead of a learned matrix, and also lacks ``input'' and ``output'' gates, so for a given hidden size, it requires about half the parameters.

The LMU is a version of the HiPPO-LegT cell with an additional hidden-to-hidden transition matrix and memory-to-memory transition vector instead of the gate $g$, leaving it with approximately the same number of trainable parameters.

\paragraph{Training Details}
Unless stated otherwise, all methods use the Adam optimizer~\citep{kingma2014adam} with learning rate frozen to $0.001$, which has been a robust default for RNN based models~\citep{voelker2019legendre, gu2020improving}.

All experiments use PyTorch 1.5 and are run on a Nvidia P100 GPU.









\subsection{Permuted MNIST}
\label{sec:pmnist-details}

\paragraph{Task}
The input to the sequential MNIST (sMNIST) task \cite{le2015simple} is an MNIST source image,
flattened in row-major order into a single sequence of length 784.
The goal of the model is to process the entire image sequentially before
outputting a classification label, requiring learning long-term dependencies.
A variant of this, the permuted MNIST (pMNIST) task, applies a fixed permutation
to every image, breaking locality and further straining a model's capacity for
long-term dependencies.

Models are trained using the cross-entropy loss.
We use the standard train-test split (60,000 examples for training and 10,000
for testing), and further split the training set with
10\% to be used as validation set.

\paragraph{Baselines and Ablations}

\cref{tab:pmnist} is duplicated here in \cref{tab:pmnist-val,tab:pmnist-test}, with more complete baselines and hyperparameter ablations.

\cref{tab:pmnist-val} consists of our implementations of various baselines related to our method, described in \cref{subsec:model_architectures}.
Each method was ran for 3 seeds, and the maximum average validation accuracy is reported.

All methods used the same hidden size of $512$; we found that this gave better performance than $256$, and further increasing it did not improve more.
All methods were trained for 50 epochs with a batch size of 100.

\paragraph{State of the Art}
\cref{tab:pmnist-test} directly shows the reported test accuracy of various methods on this data (Middle and Bottom).
\cref{tab:pmnist-test} (Top) reports the test accuracy of various instantations of our methods.
We additionally include our reproduction of the LMU,
which achieved better results than reported in~\citet{voelker2019legendre} (possibly due to a larger hidden size).
We note that \emph{all} of our HiPPO methods are competitive; each of them (HiPPO-LegT, HiPPO-LagT, HiPPO-LegS) achieves state-of-the-art among previous recurrent sequence models.
Note that differences between our HiPPO-LegT and LMU numbers in \cref{tab:pmnist-test} (Top) stem primarily from the architecture difference (\cref{subsec:model_architectures}).

\paragraph{Timescale Hyperparameters}
\cref{tab:pmnist-val} also shows ablations for the HiPPO-LegT and HiPPO-LagT timescale hyperparameters.
HiPPO-LagT sweeps the discretization step size $\Delta t$ (\cref{sec:discretization,sec:discretization-full}).
For LegT, we set $\Delta t = 1.0$ without loss of generality, as only the ratio of $\theta$ to $\Delta t$ matters.
These timescale hyperparameters are important for these methods.
Previous works have shown that the equivalent of $\Delta t$ in standard RNNs, i.e.\ the gates of LSTMs and GRUs (\cref{sec:discretization}), can also drastically affect their performance \cite{tallec2018can,gu2020improving}.
For example, the only difference between the URLSTM and LSTM in \cref{tab:pmnist-test} is a reparametrization of the gates.







\begin{table}[ht]%
  \caption{
    \textbf{Our methods and related baselines}. Permuted MNIST (pMNIST) validation scores. (Top): Our methods. (Bottom): Recurrent baselines.
  }
    \centering
    \begin{tabular}{@{}ll@{}}
      \toprule
      Method                     & Validation accuracy (\%) \\
      \midrule
      \textbf{HiPPO-LegS}        & \textbf{98.34}           \\
      HiPPO-LagT $\Delta t = 1.0$           & 98.15                    \\
      HiPPO-LegT $\theta = 200$  & 98.00                     \\
      HiPPO-LegT $\theta = 2000$ & 97.90                     \\
      HiPPO-LagT $\Delta t = 0.1$           & 96.44                    \\
      HiPPO-LegT $\theta = 20$   & 91.75                    \\
      HiPPO-LagT $\Delta t = 0.01$           & 90.71                    \\
      HiPPO-Rand                 & 69.93 \\
      \midrule
      LMU                        & 97.08                    \\
      ExpRNN                     & 94.67                    \\
      GRU                        & 93.04                    \\
      LSTM                       & 92.54                    \\
      MGU                        & 89.37                    \\
      RNN                        & 52.98                    \\
      \bottomrule
    \end{tabular}
    \label{tab:pmnist-val}
  \end{table}
\begin{table}[ht]
  \caption{ \textbf{Comparison to prior methods for pixel-by-pixel image classification.} Reported test accuracies from previous works on pixel-by-pixel image classification benchmarks. Top: Our methods. Middle: Recurrent baselines and variants. Bottom: Non-recurrent sequence models with global receptive field.
  }
    \centering
    \begin{tabular}{@{}lll@{}}
      \toprule
      Model                                     & Test accuracy (\%) \\
      \midrule
      \textbf{HiPPO-LegS}                       & \textbf{98.3}      \\
      HiPPO-Laguerre                            & 98.24              \\
      HiPPO-LegT                                & 98.03              \\
      LMU (ours)                                & 97.29              \\
      \midrule
      URLSTM + Zoneout \citep{krueger2016zoneout}          & 97.58              \\
      LMU~\citep{voelker2019legendre}           & 97.15              \\
      URLSTM~\citep{gu2020improving}            & 96.96              \\
      IndRNN~\citep{indrnn}                     & 96.0               \\
      Dilated RNN~\citep{chang2017dilated}      & 96.1               \\
      r-LSTM ~\citep{trinh2018learning}         & 95.2               \\
      LSTM~\citep{gu2020improving}              & 95.11              \\
      \midrule
      TrellisNet~\citep{trellisnet}             & 98.13              \\
      Temporal ConvNet~\citep{bai2018empirical} & 97.2               \\
      Transformer~\citep{trinh2018learning}     & 97.9               \\
      \bottomrule
    \end{tabular}
    \label{tab:pmnist-test}
\end{table}


\subsection{Copying}
\label{sec:copying-details}

\paragraph{Task}
In the Copying task~\citep{arjovsky2016unitary},
the input is a sequence of $L+20$ digits where the first 10 tokens $\left(a_0, a_1, \dots, a_9\right)$ are randomly chosen from $\left\{1,\dots,8 \right\}$, the middle N tokens are set to $0$, and the last ten tokens are $9$.
The goal of the recurrent model is to output $\left(a_0, \dots, a_9\right)$ in order on the last 10 time steps, whenever the cue token $9$ is presented.
Models are trained using the cross-entropy loss; the random guessing baseline has loss $\log(8) \approx 2.08$.
We use length $L=200$.
The training and testing examples are generated in the same way.

Our motivation of studying the Copying task is that standard models such as the LSTM struggle to solve it.
We note that the Copying task is much harder than other memory benchmarks such as the Adding task~\cite{arjovsky2016unitary}, and we do not consider those.

\begin{figure}[ht]
  \centering
  \includegraphics[width=.7\linewidth]{figures/copying200.png}
  \caption{Loss on the Copying task.
    HiPPO methods are the only to fully solve the task.
    The hyperparameter-free LegS update is best, while methods with timescale parameters (e.g.\ LegT) do not solve the task if mis-specified.
  }
  \label{fig:copy200}
\end{figure}

\paragraph{Results}
The HiPPO-LegS method solves this task the fastest.
The LegT method also solves this task quickly, only if the parameter $\theta$ is initialized to the correct value of $200$.
Mis-specifying this timescale hyperparameter to $\theta=20$ or $\theta=2000$ drastically slows down the convergence of HiPPO-LegT.
The LMU (at optimal parameter $\theta=200$) solves this task at comparable speed;
like in \cref{sec:pmnist-details}, differences between HiPPO-LegT $(\theta=200)$ and LMU here arise from the minor architecture difference in \cref{subsec:model_architectures}.

The HiPPO-Rand baseline (denoted ``random LTI'' system here)
does much worse than the updates with the dynamics derived from our framework,
highlighting the importance of the precise dynamics (in contrast to just the architecture).

Standard methods such as the RNN and LSTM are also nearly stuck at baseline.


\subsection{Trajectory Classification}
\label{sec:trajectory-details}

\paragraph{Dataset}
The Character Trajectories dataset~\citep{bagnall2018uea} from the UCI machine learning repository~\citep{Dua:2019} consists of pen tip trajectories recorded from writing individual characters.
The trajectories were captured at 200Hz and data was normalized and smoothed.
Input is 3-dimensional ($x$ and $y$ positions, and pen tip force),
and there are 20 possible outputs (number of classes).
Models are trained using the cross-entropy loss.
The dataset contains 2858 time series.
The length of the sequences is variable, ranging up to $182$.
We use a train-val-test split of 70\%-15\%-15\%.

\paragraph{Methods}

RNN baselines include the LSTM~\citep{lstm}, GRU~\citep{chung2014empirical}, and LMU~\citep{voelker2019legendre}.
Our implementations of these used 256 hidden units each.

The GRU-D~\citep{che2018recurrent} is a method for handling missing values in time series that computes a decay between observations.
The ODE-RNN~\citep{rubanova2019latent} and Neural CDE (NCDE)~\citep{kidger2020neural} baselines are state-of-the-art neural ODE methods, also designed to handle irregularly-sampled time series.
Our GRU-D, ODE-RNN, and Neural CDE baselines used code from~\citet{kidger2020neural}, inheriting the hyperparameters for those methods.


All methods trained for 100 epochs.

\paragraph{Timescale mis-specification}

The goal of this experiment is to investigate the performance of models when the timescale is mis-specified between train and evaluation time, leading to distribution shift.
We considered the following two standard types of time series:
\begin{enumerate}%
    \item Sequences sampled at a fixed rate
    \item Irregularly-sampled time series (i.e., missing values) with timestamps
\end{enumerate}

Timescale shift is emulated in the corresponding ways, which can be interpreted as different sampling rates or trajectory speeds.
\begin{enumerate}%
    \item Either the train or evaluation sequences are downsampled by a factor of 2
    \item The train or evaluation timestamps are halved.\footnote{Instead of the train timestamps being halved, equivalently the evaluation timestamps can be doubled.}
\end{enumerate}
The first scenario in each corresponds to the original sequence being sampled at 100Hz instead of 200Hz; alternatively, it is equivalent to the writer drawing twice as fast.
Thus, these scenarios correspond to a train $\to$ evaluation timescale shift of 100Hz $\to$ 200Hz and 200Hz $\to$ 100Hz respectively.

Note that models are unable to obviously tell that there is timescale shift.
For example, in the first scenario, shorter or longer sequences can be attributed to the variability of sequence lengths in the original dataset.
In the second scenario, the timestamps have different distributions, but this can correspond to different rates of missing data, which the baselines for irregularly-sampled data are able to address.






\subsection{Online Function Approximation and Speed Benchmark}
\label{subsec:function_approx}

\paragraph{Task} The task is to reconstruct an input function (as a discrete
sequence) based on some hidden state produced after the model has traversed the
input function.
This is the same problem setup as in Section~\ref{subsec:hippo-setup};
the online approximation and reconstruction details are in \cref{sec:hippo-framework-details}.
The input function is randomly sampled from a continuous-time band-limited white
noise process, with length $10^6$.
The sampling step size is $\Delta t = 10^{-4}$, and the signal band limit is 1Hz.

\paragraph{Models} We compare HiPPO-LegS, LMU, and LSTM.
The HiPPO-LegS and LMU model only consists of the memory update and not the
additional RNN architecture.
The function is reconstructed from the coefficients using the formula in
\cref{sec:derivations}, so no training is required.
For LSTM, we use a linear decoder to reconstruct the function from the LSTM
hidden states and cell states, trained on a collection of 100 sequences.
All models use $N = 256$ hidden units.
The LSTM uses the $L2$ loss.
The HiPPO methods including LMU follow the fixed dynamics of \cref{thm:legt-lagt} and \cref{thm:legs}.

\paragraph{Speed benchmark}
We measure the inference time of HiPPO-LegS, LMU, and LSTM, in
single-threaded mode on a server Intel Xeon CPU E5-2690 v4 at 2.60GHz.


\subsection{Sentiment Classification on the IMDB Movie Review Dataset}
\label{sec:imdb}

\paragraph{Dataset} The IMDB movie review
dataset~\citep{maas-EtAl:2011:ACL-HLT2011} is a standard binary sentiment
classification task containing 25000 train and test sequences, with sequence
lengths ranging from hundreds to thousands of steps.
The task is to classify the sentiment of each movie review into either positive
or negative.
We use 10\% of the standard training set as validation set.

\paragraph{Methods}
RNN baselines include the LSTM~\citep{lstm}, vanilla RNN,
LMU~\citep{voelker2019legendre}, and expRNN~\citep{lezcano2019cheap}.
Our implementations of these used 256 hidden units each.

\paragraph{Result}
As shown in Table~\ref{tab:imdb_test_acc}, our HiPPO-RNNs have similar and consistent performance,
on par or better than LSTM.
Other long-range memory RNN approaches that constrains the expressivity of the
network (e.g.\ expRNN) performs worse on this more generic task.


\begin{table}
    \centering
    \begin{tabular}{@{}ll@{}}
      \toprule
      Model                  & Test accuracy (\%) \\
      \midrule
      HiPPO-LegS             & 87.8 $\pm$ 0.2 \\
      HiPPO-LagT         & \textbf{88.0} $\pm$ 0.2 \\
      HiPPO-LegT $\theta = 100$   & 87.4 $\pm$ 0.3 \\
      HiPPO-LegT $\theta = 1000$  & 87.7 $\pm$ 0.2 \\
      HiPPO-LegT $\theta = 10000$ & 87.9 $\pm$ 0.3 \\
      HiPPO-Rand             & 82.9 $\pm$ 0.3 \\ \midrule
      LMU $\theta = 1000$         & 87.7 $\pm$ 0.1 \\
      LSTM                   & 87.3 $\pm$ 0.4 \\
      expRNN                 & 84.3 $\pm$ 0.3 \\
      RNN                    & 67.4 $\pm$ 7.7 \\
      \bottomrule
    \end{tabular}
    \caption{IMDB test accuracy, averaged over 3 seeds. Top: Our methods.
      Bottom: Recurrent baselines.}
    \label{tab:imdb_test_acc}
\end{table}

\subsection{Mackey Glass prediction}
\label{sec:mackey}


The Mackey-Glass data~\citep{mackey1977oscillation}
is a time series prediction task for modeling chaotic dynamical systems.
We build on the implementation of~\citet{voelker2019legendre}.
The data is a sequence of one-dimensional observations,
and models are tasked with predicting 15 time steps into the future.
The models are 4-layer stacked recurrent neural networks, trained with the mean squared error (MSE) loss.
\citet{voelker2019legendre} additionally consider a hybrid model with alternating LSTM and LMU layers,
which improved on either by itself.
We did not try this approach with our method HiPPO-LegS such as combining it with the LSTM or other HiPPO methods,
but such ideas could further improve our performance.
As a baseline method, the identity function does not simulate the dynamics,
and simply guesses that the future time step is equal to the current input.


\cref{fig:mackey} plots the training and validation mean squared errors (MSE) of these methods.
The table reports final normalized root mean squared errors (NRMSE) $\sqrt{\frac{\mathbb{E}\left[ (Y-\hat{Y})^2 \right]}{\mathbb{E}[Y^2]}}$ between the targets $Y$ and predictions $\hat{Y}$.
HiPPO-LegS outperforms the LSTM, LMU, and the best hybrid LSTM+LMU
model from [68], reducing normalized MSE by over 30\%.


\begin{figure}[!ht]
  \begin{minipage}{.5\linewidth}%
    \centering
    \includegraphics[width=\linewidth]{figures/mackey.pdf}
  \end{minipage}
  \begin{minipage}{.5\linewidth}%
    \begin{tabular}{@{}lll@{}}
      \toprule
      Model    & Test MSE              & Test NRMSE \\
      \midrule
      Baseline & 0.1229                & 1.62274    \\
      LSTM     & 4.784e-4              & 0.10123    \\
      LMU      & 4.414e-4              & 0.09722    \\
      Hybrid LSTM+LMU  & 2.198e-4      & 0.06862    \\
      LegS     & 1.054e-4              & 0.04752    \\
      \bottomrule
    \end{tabular}
  \end{minipage}
  \caption{Mackey-Glass predictions}
  \label{fig:mackey}
\end{figure}


\subsection{Additional Analysis and Ablations of HiPPO}
\label{sec:exp-hippo-ablations}

To further analyze the tradeoffs of the memory updates derived from our framework,
in \cref{fig:function_approx_full} we plot a simple input function
$f(x) = 1/4 \sin x + 1/2 \sin (x/3) + \sin(x/7)$ to be approximated.
The function is subsampled on the range $x \in [0, 100]$, creating a sequence of length $1000$.
This function is simpler than the functions sampled from white noise
signals described in \cref{subsec:function_approx}.
Given this function, we use the same methodology as in \cref{subsec:function_approx} for processing the function online and then reconstructing it at the end.

In Figure~\ref{fig:function_approx_full}(a, b), we plot the true function $f$, and
its absolute approximation error based on LegT, LagT, and LegS.
LegS has the lowest approximation error, while LegT and LagT are similar and slightly worse than LegS.
Next, we analyze some qualitative behaviors.

\paragraph{LegT Window Length}

In Figure~\ref{fig:function_approx_full}(c), shows that the approximation error of
LegT is sensitive to the hyperparameter $\theta$, the length of the window.
Specifying $\theta$ to be even slightly too small (by $0.5\%$ relative to the total sequence length)
causes huge errors in approximation.
This is expected by the HiPPO framework, as the final measure $\mu^{(t)}$ is not supported everywhere, so the projection problem does not care that the reconstructed function is highly inaccurate near $x=0$.

\paragraph{Generalized LagT Family}
Our LagT method actually comprises a family of related transforms,
governed by two parameters $\alpha, \beta$ specifying the original measure and the tilting (\cref{sec:derivation-lagt}).
\cref{fig:function-approx-lagt} shows the error as these parameters change.
\cref{fig:function-approx-lagt}(a) shows that small $\alpha$ generally performs better.
\cref{fig:function-approx-lagt}(b, c) show that the reconstruction is unstable for larger $\beta$, but small values of $\beta$ work well.
More detailed theoretical analysis explaining these tradeoffs would be an interesting question to analyze.

\paragraph{LegS vs. LegT}

In comparison to LegT, LegS does not need any hyperparameters governing the timescale.
However, suppose that the LegT $\theta$ window size was chosen perfectly to match the length of the sequence;
that is, $\theta = T$ where $T$ is the final time range.
Note that at the end of consuming the input function (time $t = T$),
the measures $\mu^{(t)}$ for LegS and LegT are \emph{both} equal to $\frac{1}{T}\mathbb{I}_{[0,T]}$ (\cref{subsec:high_order_projection,sec:theory_legs}).
Therefore, the approximation $\proj_T(f)$ is specifying the same function for both LegS and LegT at time $t=T$.
The sole difference is that LegT has an additional approximation term for $f(t-\theta)$ while calculating the update at every time $t$ (see \cref{sec:derivation-legt}), due to the nature of the sliding rather than scaling window.


\begin{figure}[ht]
  \centering
  \begin{subfigure}[t]{0.33\linewidth}
    \centering
    \includegraphics[width=\textwidth]{figures/function_approx_true_f}
    \caption{True function $f(x)$}
  \end{subfigure}%
  \begin{subfigure}[t]{0.33\linewidth}
    \centering
    \includegraphics[width=\textwidth]{figures/function_approx_error_bilinear}
    \caption{Absolute approx.\ error}
  \end{subfigure}%
  \begin{subfigure}[t]{0.33\linewidth}
    \centering
    \includegraphics[width=\textwidth]{figures/function_approx_error_theta}
    \caption{Error for different $\theta$'s in LegT}
  \end{subfigure}%
  \caption{Function approximation comparison between LegT, LagT, and LegS. LegS
  has the lowest approximation error. LegT error is sensitive to the choice of
  window length $\theta$, especially if $\theta$ is smaller than the length of the true function.}
  \label{fig:function_approx_full}
\end{figure}

\begin{figure}[ht]
  \centering
  \begin{subfigure}[t]{0.33\linewidth}
    \centering
    \includegraphics[width=\textwidth]{figures/laguerre_alpha.pdf}
    \caption{Generalized Laguerre family,\\ fixed $\beta=0.01$ and varying $\alpha$}
  \end{subfigure}
  \begin{subfigure}[t]{0.33\linewidth}
    \centering
    \includegraphics[width=\textwidth]{figures/laguerre_small_beta.pdf}
    \caption{Generalized Laguerre family,\\ fixed $\alpha=0$ and small $\beta$}
  \end{subfigure}%
  \begin{subfigure}[t]{0.33\linewidth}
    \centering
    \includegraphics[width=\textwidth]{figures/laguerre_large_beta.pdf}
    \caption{Generalized Laguerre family,\\ fixed $\alpha=0$ and large $\beta$}
  \end{subfigure}%
  \caption{
    Function approximation comparison between different instantiations of the generalized tilted Laguerre family (\cref{sec:derivation-lagt}).
  }
  \label{fig:function-approx-lagt}
\end{figure}

\end{document}

